\section{Regular moments, singular ranks and recursive exponents}\label{sec:spectral-conversion}
We have completed one entropy-to-moment route. Before changing the entropy or sparse input, we record its consequences for singular ranks and repeated sweeps. We then isolate the moment amplification used by several later balanced recursions. Their entropy laws, sparse estimates and representation cutoffs remain separate inputs.

Let $G$ be a finite group and $B$ convolution by a probability law on $G$, acting on $\ell^2(G)$ with counting measure. Both $B$ and $B^*$ fix constants and preserve their orthogonal complement. Write $B_0$ for the restriction to that complement. For an irreducible unitary representation $\rho$, let $D_\rho$ be its dimension and $B_\rho$ its Fourier matrix on one copy. Its singular values appear $D_\rho$ times in the regular representation. We include zeros in every singular list, ordered decreasingly.

\begin{proposition}[Moment and rank conversion]\label{e:nonnormal-moment-conversion}
Suppose, for $p>0$ and $0<\varepsilon\le1/2$, that
\begin{equation}\label{e:nonconstant-moment}
 \sum_{\rho\ne\mathbf1}D_\rho\Tr|B_\rho|^p\le\varepsilon.
\end{equation}
Then every nontrivial $\rho$ and every $1\le r\le D_\rho$ satisfy
\begin{equation}\label{e:block-rank}
 s_r(B_\rho)\le\left(\frac{\varepsilon}{D_\rho r}\right)^{1/p}.
\end{equation}
The full regular singular list, with its constant value at index $1$, satisfies
\begin{equation}\label{e:regular-rank}
 s_j(B)\le j^{-1/p}\qquad(1\le j\le|G|).
\end{equation}
For every integer $M$ with $2M\ge p$, $M$ independent convolution steps have worst-start total-variation distance at most $\tfrac12\sqrt\varepsilon$ from uniform.

Conversely, if $s_j(B)\le j^{-c}$ for an absolute $c>0$ over a family of finite groups and their convolution operators, then for each fixed $\varepsilon>0$ there is a fixed $p'>0$, independent of the group, for which \eqref{e:nonconstant-moment} holds with $p'$ in place of $p$.
\end{proposition}
\begin{proof}
The first $r$ singular values in $B_\rho$ are at least $s_r(B_\rho)$, so their regular contribution is at least $D_\rho r\,s_r(B_\rho)^p$. This proves \eqref{e:block-rank}. It applies also when that singular value is zero.

The left side of \eqref{e:nonconstant-moment} is precisely the sum of the $p$th powers of all nonconstant regular singular values. In particular none equals $1$. Thus the distinguished constant singular value is the first value and is simple. For $j\ge2$,
\[
 (j-1)s_j(B)^p\le\varepsilon\le\frac{j-1}{j}.
\]
This proves \eqref{e:regular-rank}; at $j=1$ it is equality.

The power estimate uses the matrix product, rather than a spectral identification. Schatten H\"older with $M$ factors gives
\begin{equation}\label{e:holder-power}
 \|B_0^M\|_{\HS}^2
 \le\|B_0\|_{2M}^{2M}
 =\sum_{j\ge2}s_j(B)^{2M}
 \le\sum_{j\ge2}s_j(B)^p\le\varepsilon.
\end{equation}
The last comparison uses $s_j(B)\le1$. If $\mu_M$ is the convolution law after $M$ steps and $\Pi$ projects onto constants, every row of $B^M-\Pi$ is a permutation of $\mu_M-|G|^{-1}$. Consequently
\[
 \|B_0^M\|_{\HS}^2
 =|G|\sum_{g\in G}|\mu_M(g)-|G|^{-1}|^2.
\]
Cauchy--Schwarz bounds total variation by one half of this square root, uniformly over the starting group element.

For the converse the nonconstant regular moment is at most
\[
 \sum_{j=2}^{|G|}j^{-cp'}\le\sum_{j=2}^\infty j^{-cp'}.
\]
For $q>1$ the last series with exponent $q$ is at most
$2^{-q}+2^{1-q}/(q-1)$, by integral comparison. It tends to zero as $q\to\infty$. Choosing $q=cp'$ sufficiently large proves the uniform assertion.
\end{proof}

\begin{corollary}[Vanishing error after a fixed number of sweeps]\label{e:vanishing-mixing}
Fix $p>0$ and $0<\varepsilon\le1/2$. Suppose \eqref{e:nonconstant-moment} holds uniformly for a sequence of groups, and every nonconstant block which is not zero has dimension at least $D_{\min}\to\infty$. For any fixed integer $M$ with $2M>p$,
\[
 4\|\mu_M-U_G\|_{\TV}^2
 \le\varepsilon\left(\frac{\varepsilon}{D_{\min}}\right)^{(2M-p)/p}
 \longrightarrow0.
\]
\end{corollary}
\begin{proof}
Equation~\eqref{e:block-rank} with $r=1$ bounds
$\|B_0\|_{\op}\le(\varepsilon/D_{\min})^{1/p}$. Retain the extra powers in \eqref{e:holder-power} and bound each by this supremum:
\[
 \sum_{j\ge2}s_j(B)^{2M}
 \le \|B_0\|_{\op}^{2M-p}\sum_{j\ge2}s_j(B)^p.
\]
The preceding proof identifies this bound with an upper bound for four times total variation squared.
\end{proof}
For $S_n$, the sign block of a sweep vanishes by Lemma~\ref{lem:annihilation}. The minimum dimension among the other nonconstant irreducibles tends to infinity: otherwise their reciprocal-degree sum could not tend to zero in Lemma~\ref{lem:degrees}. Thus Theorem~\ref{tra:duality:moment} gives vanishing worst-start error after any fixed integer $M>p_*$. Its positive moment has exponent $p_*$, whereas its Schatten moment has exponent $2p_*$. This factor of two will remain explicit when later estimates use one convention or the other.

There is also a useful consequence when an argument gives only irreducible indexed bounds. If $s_r(B_\rho)\le(D_\rho r)^{-a}$, then $\|B_\rho\|_{\op}\le D_\rho^{-a}$ and
\[
 D_\rho\|B_\rho^M\|_{\HS}^2\le D_\rho^{2-2aM}.
\]
For symmetric groups, taking $2aM\ge3$, annihilating sign separately, and using Lemma~\ref{lem:degrees} makes the sum tend to zero. This deduction does not assume normality either. It compares the final absolute-power conclusions; it does not identify the constants, sparse ranges, or retained weights in the proofs that follow.

\subsection{Closing a balanced moment recursion}

\begin{lemma}[Exponent losses under rounded halving]\label{lem:dyadic-exponents}
Fix $\gamma>0$ and an integer base $d_0\ge3$. Along any branch
obtained by replacing $d>d_0$ with $\lfloor d/2\rfloor$ or
$\lceil d/2\rceil$, the sum of $d^{-\gamma}$ over the terms
above the base is uniformly bounded. For each fixed $C\ge0$, products of
$1+C d^{-\gamma}$ are uniformly bounded, and products of
$1-C d^{-\gamma}$ are bounded away from zero if each factor
is at least $1/2$.
\end{lemma}
\begin{proof}
Each child is at most $2/3$ of its parent. Read upward from the
last term above the base: the summands decrease by a factor at
most $(2/3)^\gamma$, so their sum is bounded by a geometric
series. The product assertions follow from
$\log(1+x)\le x$ and $\log(1-x)\ge-2x$ for $0\le x\le1/2$.
\end{proof}

The grid estimates below first give a large bound on a regular trace.
Such a bound is useful because a singular value in a high-dimensional
irreducible occurs many times in that trace. Increasing the moment
exponent then removes the large-dimensional contribution. The following
lemma separates this common step from the estimates specific to each
entropy exposure.

Fix $\nu\in\{1,2\}$ and write
\[
 Z_n(p)=\operatorname{Tr}_{\rm reg}|T_n|^{\nu p}.
\]
Thus $\nu=2$ uses the positive-square convention, whereas $\nu=1$
uses the singular-value convention. The parameter $p$ has this same
meaning throughout one application of the lemma.

\begin{lemma}[Moment amplification at a balanced split]
\label{lem:balanced-moment-closure}
Fix $0<\varepsilon<1$ and $P>0$. Suppose that, above an absolute
bit length $d_0\ge3$, the nonconstant irreducible blocks at $n=2^d$ are divided
into two classes with the following properties. Each class includes
all regular copies of every singular value in its blocks.
\begin{enumerate}
\item For every $p\ge P$, the first class contributes at most
$\varepsilon/2$ to $Z_n(p)$.
\item Each value in the second class occurs with its irreducible
regular multiplicity at least $e^{H_n}$.
\item Put $a=2^{\lfloor d/2\rfloor}$ and
$b=2^{\lceil d/2\rceil}$. For every $p\ge P$,
\[
 Z_a(p),Z_b(p)\le1+\varepsilon
 \quad\Longrightarrow\quad
 Z_n(\alpha_d p)\le e^{E_n},
\]
where $\alpha_d\ge1$ and the size threshold is independent of $p$.
\end{enumerate}
Suppose further that numbers $\delta_d>0$ satisfy
\[
 (1+\delta_d)E_n-\delta_d H_n\le\log(\varepsilon/2),
 \qquad
 \alpha_d(1+\delta_d)\le1+C d^{-\gamma}
\]
for some fixed $C<\infty$ and $\gamma>0$.
Then $Z_n(p_*)\le1+\varepsilon$ for all dyadic $n$, at one
finite exponent $p_*$.
\end{lemma}

\begin{proof}
Assume the child bounds at parameter $p$ and set $t=\alpha_dp$.
If $s$ belongs to the second class, regular multiplicity and the
rough trace estimate imply
\[
 s^{\nu t}\le e^{E_n-H_n}.
\]
Consequently its entire class contributes, at the increased parameter
$t(1+\delta_d)$, at most
\[
 \sum_{s\text{ in second class}}s^{\nu t(1+\delta_d)}
 \le e^{\delta_d(E_n-H_n)}Z_n(t)
 \le e^{(1+\delta_d)E_n-\delta_dH_n}
 \le\varepsilon/2.
\]
The sums here list regular multiplicities. The first class contributes
at most $\varepsilon/2$ by monotonicity and the first hypothesis.
The constant singular value contributes exactly one. This proves the
parent bound.

At the finitely many sizes $d\le d_0$, Lemma~\ref{lem:finitegap}
allows one common parameter $P_0\ge P$ for which the desired bound
holds. Set $p_d=P_0$ in this base range, and above it define
\[
 p_d=\alpha_d(1+\delta_d)
       \max\{p_{\lfloor d/2\rfloor},p_{\lceil d/2\rceil}\}.
\]
The preceding parent estimate proves $Z_{2^d}(p_d)\le1+\varepsilon$
by induction. Lemma~\ref{lem:dyadic-exponents} bounds the product
of the exponent multipliers along every branch of rounded halvings.
Thus $p_*:=\sup_d p_d$
is finite; contraction and monotonicity give the asserted bound at
this common parameter.
\end{proof}
